

\begin{lemma}[Volume of boxes]\label{lem:box-volume}
    Let $E = (a_{1}, b_{1}) \times \cdots \times (a_{n}, b_{n}) \subset \mathbb{R}^{n}$ where $a_{i} < b_{i}$. Then $E, \overline{ E } $ are Jordan measurable and 
    \begin{align*}
        \mu(E) = \mu(\overline{ E } ) = \prod_{i=1}^{n}(b_{i} - a_{i}).
    \end{align*}
\end{lemma}
\begin{proof}
    For $p \in \mathbb{N}$ and $t \in \mathbb{R}$, we define
    \begin{align*}
        \underline{t} = \max \left\{ 2^{-p}k \mid k \in \mathbb{Z} \text{ s.t. } 2^{-p}k < t \right\}, \quad \overline{ t }  = \min \left\{ 2^{-p}k \mid k \in \mathbb{Z} \text{ s.t. } 2^{-p}k > t \right\}.
    \end{align*}
    Notice that always $t - \underline{t}, \overline{ t } - t \leq 2^{-p}$, $t < \overline{ t } $ and $t > \underline{t}$. From this we get that for any $1 \leq i \leq n$,
    \begin{align*}
        [a_{i}, b_{i}] \subset [\underline{a_{i}}, \overline{ b_{i} }), \quad [\overline{ a_{i} }, \underline{b_{i}} ) \subset (a_{i}, b_{i}).
    \end{align*}
    Then it's clear that
    \begin{align*}
        \overline{ E } \subset \prod_{i=1}^{n} [\underline{a_{i}}, \overline{ b_{i} } ) = G_{p}, \quad E \supset \prod_{i=1}^{n} [ \overline{ a_{i} }, \underline{b_{i}} ) = F_{p}
    \end{align*}
    and $G_{p}, F_{p}$ are dyadic sets. Notice that also
    \begin{align*}
        \underline{b_{i}} - \overline{ a_{i} }  \geq b_{i} - a_{i}, \quad \overline{ b_{i} } - \underline{a_{i}} \leq b_{i} + 2^{-p} - (a_{i} - 2^{-p}) = b_{i} - a_{i} + 2^{1-p}. 
    \end{align*}
    Hence\footnote{If you find a dyadic decomposition and compute the volume, everybody knows this is true, we just had to get the boxes to have coordinates that are powers of two so it's easy to fit dyadic cubes into it.}, we get for sufficiently big $p$
    \begin{align*}
        \mu(F_{p}) = \prod_{i=1}^{n} ( \underline{ b_{i} } - \overline{ a_{i} }  ) \geq \prod_{i=1}^{n} (b_{i} - a_{i}), \quad \mu(G_{p}) = \prod_{i=1}^{n} ( \overline{ b_{i} } - \underline{ a_{i} }  ) \leq \prod_{i=1}^{n} (b_{i} - a_{i} + 2^{1-p}).
    \end{align*}
    Letting $p \to \infty$, using $\mu(F_{p}) \leq \mu_{\text{in}}(E) \leq \mu_{\text{out}}(\overline{ E } ) \leq \mu(G_{p})$, we get the statement. 
\end{proof}

\begin{prop}[Unions, intersections, differences of Jordan measurable sets]\label{prop:diff-int-un-jordan}
    Let $E_{1}$ and $E_{2}$ be Jordan measurable. Then $E_{1} \cup E_{2}$, $E_{1} \cap E_{2}$, $E_{1} \setminus E_{2}$ are all Jordan measurable and $\mu$ is additive, so if $E_{1}, E_{2}$ are Jordan measurable and disjoint then $\mu(E_{1} + E_{2}) = \mu(E_{1}) + \mu(E_{2})$.
\end{prop}
\begin{proof}
    First we prove additivity and measurability of the disjoint union. So assume that $E_{1}, E_{2}$ are disjoint. By superadditivity and subadditivity \ref{lem:prop-jordan-measure},
    \begin{align*}
        \mu_{\text{in}}(E_{1}) + \mu_{\text{in}}(E_{2}) \leq \mu_{\text{in}}(E_{1} \cup E_{2}) \leq \mu_{\text{out}}(E_{1} \cup E_{2}) \leq \mu_{\text{out}}(E_{1}) + \mu_{\text{out}}(E_{2}).
    \end{align*}
    But since the last and first terms equal, the inequality collapses to an equality.

    Now we inspect the difference. By assumption, for any $\epsilon > 0$, $\exists G_{i}, F_{i}$ dyadic such that $F_{i} \subset E_{i} \subset G_{i}$ and
    \begin{align*}
        \mu(G_{i}) - \mu(F_{i}) < \frac{\epsilon}{2}. 
    \end{align*}
    By inclusion,
    \begin{align*}
        F_{1} \setminus G_{2} \subset  E_{1} \setminus E_{2} \subset G_{1} \setminus F_{2}
    \end{align*}
    and we know that $F_{1} \setminus G_{2}$ and $G_{1} \setminus F_{2}$ are dyadic. Purely set theoretically,
    \begin{align*}
        (G_{1} \setminus F_{2}) \setminus (F_{1} \setminus G_{2}) \subset (G_{1} \setminus F_{1}) \cup (G_{2} \setminus F_{2}).
    \end{align*}
    Hence by subadditivity \ref{lem:prop-jordan-measure},
    \begin{align*}
        \mu( (G_{1} \setminus F_{2}) \setminus (F_{1} \setminus G_{2}) ) \leq \mu(G_{1} \setminus F_{1}) + \mu(G_{2} \setminus F_{2}) \leq \frac{\epsilon}{2} + \frac{\epsilon}{2} = \epsilon.
    \end{align*}

    Now we look at the intersection. Set theoretically, one writes
    \begin{align*}
        E_{1} \cap E_{2} = E_{2} \setminus (E_{2} \setminus E_{1})
    \end{align*}
    and so $E_{1} \cap E_{2}$ is Jordan measurable. And finally, we express 
    \begin{align*}
        E_{1} \cup E_{2} = (E_{1} \setminus E_{2}) \cup (E_{2} \setminus E_{1}) \cup (E_{1} \cap E_{2}),
    \end{align*}
    where the second union is disjoint, so the general union is Jordan measurable as well.
\end{proof}

\begin{lemma}[Sandwich]\label{lem:sandwich}
    Let $E \subset \mathbb{R}^{n}$. Assume for any $\epsilon > 0$, there exist $F, G \subset \mathbb{R}^{n}$ Jordan measurable such that $F \subset E \subset G$ and $\mu(G) - \mu(F) < \epsilon$. Then $E$ is Jordan measurable. 
\end{lemma}
\begin{proof}
    For any $\epsilon > 0$, there exist Jordan $F_{\epsilon} \subset E \subset G_{\epsilon}$ such that
    \begin{align*}
        \mu(F_{\epsilon}) = \mu_{\text{in}}(F_{\epsilon}) \leq \mu_{\text{in}}(E) \leq \mu_{\text{out}}(E) \leq \mu_{\text{out}}(G_{\epsilon}) = \mu(G_{\epsilon}).
    \end{align*}
    This means in particular that $\mu_{\text{out}}(E) - \mu_{\text{in}}(E) \leq \mu(G_{\epsilon}) - \mu(F_{\epsilon}) < \epsilon$. This forces $\mu_{\text{out}}(E) = \mu_{\text{in}}(E)$.
\end{proof}

\begin{definition}[Jordan null set]\label{def:jordan-null-set}
    A set $E \subset \mathbb{R}^{n}$ is called null if $\mu_{\text{out}}(E) = 0$.
\end{definition}

\begin{lemma}[Closure, interior and the Jordan measure]\label{lem:jord-clos-int}
    Let $E \subset \mathbb{R}^{n}$ be Jordan measurable. Then $E^{\circ }, \overline{ E }$ are both Jordan measurable and $\mu(E^{\circ }) = \mu(E) = \mu(\overline{ E } )$ give all the same volume.
\end{lemma}
\begin{proof}
    Let $\epsilon > 0$. Since $E$ is Jordan, $\exists F_{\epsilon}, G_{\epsilon}$ dyadic such that $F_{\epsilon} \subset E \subset G_{\epsilon}$ and $\mu(G_{\epsilon}) - \mu(F_{\epsilon}) < \epsilon$. We have
    \begin{align*}
        F_{\epsilon}^{\circ } \subset E^{\circ } \subset E \subset \overline{ E }  \subset \overline{ G }_{\epsilon}.
    \end{align*}
    Since $F_{\epsilon}$ is dyadic, it can be written as a disjoint union of Minecraft cubes
    \begin{align*}
        F_{\epsilon} = \bigcup_{\ell = 1}^{N} (a_{\ell} + [0,1)^{n})2^{-p} \implies F_{\epsilon}^{\circ } \subset \bigcup_{\ell=1}^{N} (a_{\ell} + (0,1)^{n})2^{-p}.
    \end{align*}
    And we proved that the volume of the open cube is the same as the volume of the closed cube in \ref{lem:box-volume}. Hence, for the dyadic $\mu(F_{\epsilon}^{\circ }) = \mu(F_{\epsilon})$ and the same holds for $\mu(\overline{ G }_{\epsilon}) = \mu(G_{\epsilon})$. Hence $\mu(\overline{ G }_{\epsilon}) - \mu(F_{\epsilon}^{\circ }) < \epsilon$ and passing this into the Sandwich, we get the Lemma.
\end{proof}

\begin{theorem}[Jordan measurability and boundary]\label{thm:jordan-theorem}
    $E \subset \mathbb{R}^{n}$ is Jordan measurable if and only if $E$ is bounded and $\partial E$ is null.
\end{theorem}
\begin{proof}
    $\implies$: If $E$ is Jordan measurable, then $E$ is clearly bounded. If it was unbounded, you'd need infinitely many cubes to approximate its volume from the outside. Also, by Lemma \ref{lem:jord-clos-int}, 
    \begin{align*}
        \mu(\overline{ E } ) = \mu(E^{\circ }) = \mu(E).
    \end{align*}
    By \ref{prop:diff-int-un-jordan}, $\partial E = \overline{ E } \setminus E^{\circ }$ is Jordan measurable and by additivity (the interior and boundary are disjoint)
    \begin{align*}
        \mu(\overline{ E } ) = \mu(E^{\circ }) + \mu( \partial E) \implies \mu( \partial E ) = 0.
    \end{align*}

    $\impliedby$: Assume that $E$ is bounded and $\mu_{\text{out}}(\partial E ) = 0$. Given $p \in \mathbb{N}$, let 
    \begin{align*}
        F_{p} = \bigcup \left\{ Q \text{ dyadic cube of size } 2^{-p} \mid Q \subset E^{\circ } \right\}
    \end{align*}
    and
    \begin{align*}
        G_{p} = \bigcup \left\{ Q \text{ dyadic cube of size } 2^{-p} \mid Q \cap \overline{ E } \neq \emptyset \right\}.
    \end{align*}
    Then we clearly have $F_{p} \subset E^{\circ } \subset \overline{ E } \subset G_{p}$. Since $\partial E$ is null, it is Jordan measurable and for $\epsilon > 0$, $\exists H_{\epsilon}$ such that $\partial E \subset H_{\epsilon}$ is dyadic and $\mu(H_{\epsilon}) < \epsilon$.

    Let $p$ be the pixel size of $H_{\epsilon}$. Then
    \begin{align*}
        G_{p} \setminus F_{p} \subset H_{\epsilon} \implies \mu(G_{p}) - \mu(F_{p}) < \epsilon.
    \end{align*}
    Hence $E$ is Jordan measurable by \ref{lem:sandwich}. And so it remains to justify the inclusion. Indeed, if $z \in G_{p} \setminus F_{p}$, then there exists a cube of size $2^{-p}$, call it $Q_{p}$, such that $z \in Q_{p}$. But this $Q_{p}$ must satisfy $Q_{p} \not \subset E^{\circ }$ and $Q_{p} \cap \overline{ E } \neq \emptyset$ at the same time. Hence there exist points $x,y \in Q_{p}$ such that $y \in \mathbb{R}^{n} \setminus E^{\circ }$ and $x \in \overline{ E } $. Because $Q_{p}$ is convex, we can draw a line between $x$ and $y$ and define
    \begin{align*}
        t_{*} = \sup \left\{ t \in (0,1) \mid (1-t)x+ty \in E  \right\}.
    \end{align*}
    % image here would be so fucking golden, show really well how the point is approached both from the closure and strict exterior, hence has to lie on the boundary
    But this means that the point $z_{*} = (1-t_{*})x + t_{*}y$ can be approached by sequences both from $\mathbb{R}^{n} \setminus E \subset \mathbb{R}^{n} \setminus E^{\circ }$ and from $E \subset \overline{ E } $. But since both of those sets are closed, $z_{*}$ must lie in the intersection, so $\partial E$. But that means that $Q_{p}$ has non-empty intersection with the dyadic set $H_{\epsilon}$ of pixel size $2^{-p}$, hence is entirely contained in it (easy to check), so $z \in H_{\epsilon}$.
\end{proof}

